El ordenamiento de arreglos unidimensionales y la multiplicación de matrices cuadradas constituyen dos de los problemas computacionales más estudiados en la ciencia de la computación, sirviendo como arquetipos para el análisis de técnicas como divide y vencerás, partición y optimización de bucles. Si bien el análisis asintótico tradicional ($\mathcal{O}, \Omega, \Theta, o, \omega$) describe el crecimiento del número de operaciones en el límite cuando el tamaño de la entrada tiende a infinito, en la práctica el rendimiento de una implementación se ve fuertemente influenciado por factores constantes, la arquitectura del procesador y la jerarquía de memoria caché. El objetivo de este trabajo es contrastar experimentalmente las cotas teóricas de cuatro algoritmos de ordenamiento (\textit{std::sort}, \textit{Merge Sort}, \textit{Quick Sort} y \textit{Patience Sort}) y dos de multiplicación matricial (\textit{Naive} y \textit{Strassen}), analizando cómo la estructura de los datos, el tamaño del dominio y el diseño del algoritmo determinan su comportamiento empírico en tiempo de ejecución y uso de memoria.